\section{Conclusion}
\label{sec:conclusion}

We run a controlled study of verifiable instructions in system prompts across 8 open-weight models. Instruction count is the main driver of non-compliance: per-instruction compliance falls by about 1.5--4\pp per added instruction for capable models, and the all-satisfied rate falls to 37.0\% at five instructions and 5.0\% at ten. Placing an instruction last (+1.7\pp) and repeating it (+3.3\pp) help only a little, and repetition has no cost beyond prompt length. Interference concentrates in specific pairs, especially instructions that fix the whole response layout, rather than in broad format-versus-persona competition. Smaller models do worse at every load, but they are not clearly less robust to added load once floor and ceiling effects are considered.

The key takeaway is simple: to improve compliance, remove instructions rather than rearrange them.

Next, we plan to test frontier models with $k$ up to 20--30, use judged persona and tone instructions, test output templates as a mitigation for layout capture, sample at $T>0$ with multiple seeds to separate decoding noise from small effects, and test whether pairwise interference in E2 predicts failures at higher $k$.
